%==========================================================
% ĐÁP ÁN CHI TIẾT — BÀI TẬP VỀ NHÀ
% Bài 6: NHIỆT NÓNG CHẢY RIÊNG — Vật lý 12
% Biên soạn: Kha Khung Hiệp — Đòn Bẩy AI
%==========================================================

\documentclass[12pt,a4paper]{article}
\usepackage[utf8]{vietnam}
\usepackage{amsmath, amssymb, amsfonts}
\usepackage{geometry}
\usepackage{booktabs}
\usepackage{array}
\usepackage{longtable}
\usepackage{enumitem}
\usepackage{xcolor}

\geometry{margin=2cm}
\setlength{\parindent}{0pt}
\setlength{\parskip}{6pt}

\definecolor{navy}{HTML}{1E3A5F}
\definecolor{green}{HTML}{4CAF50}
\definecolor{orange}{HTML}{FF9800}

\newcommand{\dapan}[1]{\textcolor{green}{\textbf{#1}}}
\newcommand{\dapansai}[1]{\textcolor{red}{\textbf{#1}}}
\newcommand{\keyword}[1]{\textcolor{navy}{\textbf{#1}}}

\begin{document}

\begin{center}
    {\LARGE \textbf{ĐÁP ÁN BÀI TẬP VỀ NHÀ}}\\[4pt]
    {\Large \textbf{Bài 6: NHIỆT NÓNG CHẢY RIÊNG}}\\[4pt]
    {\normalsize Vật lý 12 — GDPT 2018 (Kết nối tri thức)}\\[4pt]
    {\small Biên soạn: \textbf{Kha Khung Hiệp} — Đòn Bẩy AI}
\end{center}

\vspace{1em}
\hrule
\vspace{1em}

%==========================================================
% PHẦN I: TRẮC NGHIỆM (18 CÂU)
%==========================================================

\section*{PHẦN I: TRẮC NGHIỆM (18 CÂU)}

\subsection*{A. Nhận biết (6 câu)}

\begin{longtable}{>{\bfseries}cl}
\toprule
Câu & Đáp án \\ \midrule
1 & \dapan{A} \\ 
2 & \dapan{B} \\ 
3 & \dapan{B} \\ 
4 & \dapan{B} \\ 
5 & \dapan{C} \\ 
6 & \dapan{B} \\ \bottomrule
\end{longtable}

\noindent\textbf{Lời giải gợi ý:}
\begin{itemize}
    \item \textbf{C1:} Định nghĩa $\lambda$: nhiệt lượng cần làm nóng chảy 1 kg chất ở $T_{nc}$.
    \item \textbf{C2:} $\lambda = Q/m \Rightarrow$ đơn vị J/kg.
    \item \textbf{C3:} Công thức $Q = \lambda m$.
    \item \textbf{C4:} $T_{nc}$ nước đá ở áp suất chuẩn = 0°C = 273 K.
    \item \textbf{C5:} $\lambda$ là hằng số chất, chỉ phụ thuộc bản chất chất.
    \item \textbf{C6:} $\lambda$ (kJ/kg): Chì 23, Thiếc 59, Đồng 205, \dapan{Nhôm 398}.
\end{itemize}

\subsection*{B. Thông hiểu (6 câu)}

\begin{longtable}{>{\bfseries}cl}
\toprule
Câu & Đáp án \\ \midrule
7 & \dapan{B} \\ 
8 & \dapan{C} \\ 
9 & \dapan{B} \\ 
10 & \dapan{B} \\ 
11 & \dapan{B} \\ 
12 & \dapan{B} \\ \bottomrule
\end{longtable}

\noindent\textbf{Lời giải gợi ý:}
\begin{itemize}
    \item \textbf{C7:} Quá trình nóng chảy: $T$ không đổi $\to$ đoạn ngang ở 0°C.
    \item \textbf{C8:} Vôn kế đo hiệu điện thế, không dùng trong TN này.
    \item \textbf{C9:} Đá tan thu nhiệt của nước $\to$ nước nguội xuống.
    \item \textbf{C10:} Oát kế đo công suất $P$ và thời gian $t$ (hoặc tích hợp $Pt$).
    \item \textbf{C11:} Muối làm giảm $T_{đông đặc}$ của nước (hiện tượng giảm điểm đóng băng).
    \item \textbf{C12:} Đúc kim loại phải nấu chảy ở $T_{nc}$ của kim loại đó.
\end{itemize}

\subsection*{C. Vận dụng (6 câu)}

\begin{longtable}{>{\bfseries}cl}
\toprule
Câu & Đáp án \\ \midrule
13 & \dapan{B} \\ 
14 & \dapan{B} \\ 
\section*{PHẦN III: CÂU HỎI TRẢ LỜI NGẮN (6 CÂU)}

\begin{itemize}[label=--]
  \item \textbf{Câu 1: 680.} $Q = \lambda \cdot m = 3{,}4 \cdot 10^5 \times 2 = 680000\text{ J} = 680\text{ kJ}$.
  \item \textbf{Câu 2: 181.}
    $Q_1 = m c_{đá} \Delta t = 0{,}5 \times 2100 \times 10 = 10500\text{ J} = 10{,}5\text{ kJ}$.\\
    $Q_2 = \lambda m = 3{,}4 \cdot 10^5 \times 0{,}5 = 170000\text{ J} = 170\text{ kJ}$.\\
    $Q_{tổng} = Q_1 + Q_2 = 180{,}5\text{ kJ} \approx 181\text{ kJ}$.
  \item \textbf{Câu 3: 0.96.}
    $Q_1 = 1 \times 896 \times (658 - 20) = 571648\text{ J}$.\\
    $Q_2 = 3{,}9 \cdot 10^5 \times 1 = 390000\text{ J}$.\\
    $Q_{tổng} = 961648\text{ J} \approx 0{,}96\text{ MJ}$.
  \item \textbf{Câu 4: 7.0.}
    Nhiệt lượng tỏa hạ xuống $0^{\circ}\text{C}$ của nước ấm: $Q_{tỏa} = 200 \times 4{,}2 \times 20 = 16800\text{ J}$.\\
    Nhiệt lượng để đá tan hoàn toàn: $Q_{tan} = 30 \times 334 = 10020\text{ J}$.\\
    Dư nhiệt: $16800 - 10020 = 6780\text{ J}$.\\
    Nhiệt độ cuối: $6780 = (200 + 30) \times 4{,}2 \times t_{cb} \rightarrow t_{cb} \approx 7{,}0^{\circ}\text{C}$.
  \item \textbf{Câu 5: 5.83.}
    $Q_1 = 10 \times 380 \times (1084 - 24) = 4028000\text{ J} = 4{,}028\text{ MJ}$.\\
    $Q_2 = 10 \times 1{,}8 \cdot 10^5 = 1{,}8\text{ MJ}$.\\
    $Q_{tổng} = 4{,}028 + 1{,}8 = 5{,}828\text{ MJ} \approx 5{,}83\text{ MJ}$.
  \item \textbf{Câu 6: 2.5.}
    Nhiệt dung riêng thực nghiệm đạt $\lambda_{tn} \approx 2{,}5 \cdot 10^4\text{ J/kg}$.
\end{itemize}

\vfill
\begin{center}
\textit{--- GEMS Vật Lý 12 -- Bài 6: Nhiệt Nóng Chảy Riêng ---}
\end{center}

\end{document}